% !TeX root = ../papers/nc-p3.tex
\subsection{三生成基底の被約な点グラフの全被覆}
\label{sec:reduced-ternary-point-graphs}
ここでは三次元接点を持つ有限四辺で、基底の三関係が被約な平面三次曲線 $E$ の
自己同型 $\sigma$ のグラフを定める場合を扱う。曲線の成分を入れ替える移動も含める。

\begin{lemma}\label{lem:reduced-ternary-fourth-shift}
$\sigma$ はどの既約成分上でも点ごとの恒等にならない。
全 $E$ 上で $\sigma^4=\nu:=\mathbb P(r^{-1}\ell)$ が成立し、$\sigma$ の固定点は特異点に限る。
また $\sigma$ は既約成分の次数を保つ。
\end{lemma}
\begin{proof}
固定される成分があれば、その一般の滑らかな点 $p$ で、前節の二平面
$\operatorname{im} A_p,\operatorname{im} B_p$ は異なる。もし一致すれば
$\ell(p),r(p)$ が比例し、定理\ref{thm:elliptic-tangent-classification}の行列式の
微分計算により $p$ は特異となる。固定成分上の正規化に $L$ を引き戻した次数を $d>0$ とする。
$A_p(d_i(p,p))$ は次数 $3d$ の束の三つの切断で、その方向は一般点で
$\ell(p)\times r(p)$ の方向に入る（恒等的に零なら直ちに全点が持ち上がる）。
二次式の三切断の共通零因子を除くと、この方向への写像は次数高々 $2d$ の束から来る。
従って前者の三切断には次数少なくとも $d$ の共通零因子がある。
その零点では、グラフの仮定から特異点でも $\ker A_p=kp$ であり、
$d_i(p,p)$ が $p$ に比例する。第四座標に一次元の解が生じるので有限性に反する。
こうして各成分の一般点で $p,\sigma p$ は独立となり、前節の二平面の証明が
各成分に適用できる。被約性により $\sigma^4=\nu$ はスキーム上の等式である。
固定点での微分計算も前節通りである。

次数については前節の三次交叉を用いた曲線上の束の完全列を $n=3$ で使う。
この列の完全性には曲線の滑らかさを使わず、グラフでの階数2だけでよい。
行列式の等式は $L\otimes(\sigma^2)^*L=(\sigma^*L)^2$ を与える。
成分次数ベクトルに対する置換の作用を $P$ とすれば $(P-1)^2\deg L=0$ である。
有限位数の置換は標数零で半単純なので $P\deg L=\deg L$ となる。
\end{proof}

\begin{theorem}\label{thm:reduced-ternary-graphs}
被約三次曲線のグラフとなる上記の枝において、許される $E$ は滑らかな三次曲線、
節点三次曲線、横断する直線と滑らかな二次曲線、非共点三直線だけである。
いずれも位数2と4の二型があり、後の三種は最初の対応する型から平坦特殊化する。
横断二成分の場合は定理\ref{thm:ternary-torsion-smoothing}の $h=1$ が全てであり、
非共点三直線では定理\ref{thm:quantum-tangent-finite}の量子基底が全てである。
\end{theorem}
\begin{proof}
既約曲線は前二節による。横断する直線と二次曲線は次数により個別に保たれる。
二つの交点を交換するなら各射影直線の内部に固定点が生じるので、両交点は固定される。
座標を直線 $x=0$ と二次曲線 $yz=x^2$ とし、それぞれ $[0:1:t],[s:1:s^2]$ と表示する。
移動は $t\mapsto At$, $s\mapsto Bs$ である。$C=B^{-1}$ と置けば三関係は
\[
 R_x=yz-Azy+k x^2,\qquad R_y=zx-Cxz,\qquad R_z=xy-Cyx,\qquad k=(A-B^2)/B .
\]
両成分への代入で消え、全ての消える式がこの三本であることは
直線への制限と二次曲線上の $s^0,\ldots,s^4$ の係数から分かる。
実際 $A=B^2$ なら移動が平面に延びて平面グラフになり、他では二成分の
同時評価は階数6である（直線で階数3、その核から $s,s^2,s^3$ が出る）。
三次交叉は
\[
 xR_x+yR_y+(A/C)zR_z=R_xx+(A/C)R_yy+R_zz .
\]
従って $\nu$ は $[x:(C/A)y:(A/C)z]$ で、二成分での四乗条件は
$B^4=A B$、$A^4=(A B)^2$ となる。よって $B^4=1$, $A=C=B^{-1}$ である。
恒等を除くと $k\ne0$ で、対角変換で $k=1$ として前節の $h=1$ を得る。

非共点三直線を座標三角形とし、$\sigma$ の頂点置換を $\pi$ とする。
第0,1,2位置で順に座標を $1,\pi,\pi^2$ の順に揃えると、
各関係空間は異なる二文字の交換二項式三つからなる。
実際、頂点を揃えた各成分の移動は両端を保つ倍率であり、辺の内部で消えるためには
対角係数が零で、各非対角二項の比がこの倍率により決まる。
三次交叉のテンソルは、この座標では互いに異なる三文字の六順列だけに台を持つ。
両側で隣接交換が非零倍を保つので、繰り返し文字を隣接させた係数は零、
異なる三文字の六係数は全て非零となる。
元の座標へ戻すと台は $(i,\pi(j),\pi^2(k))$ である。
頂点 $i$ の先頭縮約 $\ell(e_i^*)$ は第一文字 $\pi(i)$ を欠く関係に、
末尾縮約 $r(e_j^*)$ は第一文字 $\pi^{-2}(j)$ を欠く関係に対応する。
従って $\nu$ の頂点置換は $\pi^3$ であり、$\pi^4=\pi^3$ より $\pi=1$ である。
これで三関係は一つの共通基底で量子交換式となり、量子基底の分類を適用できる。

接する直線と二次曲線は個別に保たれ、唯一の交点以外の各成分はアフィン直線である。
その上の移動は滑らかな固定点を持たず、いずれも非零の平行移動である。
直線 $y=0$、二次曲線 $yz=x^2$ とし、$[1:0:u]$ と $[s:1:s^2]$ で
$u\mapsto u+a,s\mapsto s+b$ と書く。
二次曲線上で $s\mapsto s+4b$ を実現する射影変換は一意で、直線上で $u\mapsto u+8b$ と作用する。
$\sigma^4$ が射影変換なので $a=2b$、従って $\sigma$ 自身が同様に射影変換となる。
被約三次曲線上で消える二次式はないので三関係は平面グラフを与え、仮定に反する。

最後に共点三直線の三重点を $p$ とする。この点は固定され、グラフの両側の評価階数は2である。
一次文字を $H=K\oplus kz$、$K=\ker p$、$p(z)=1$ とする。関係の二つの混合成分の
同時評価階数は、接空間二次元により $4-2=2$ である。
従って関係は二次元 $K$ の文字 $x_1,x_2$ を使い
\[
 R=\sum R_{ab}x_ax_b,\qquad m_j=zx_j-\sum_a F_{aj}x_a z-d_j,\quad d_j\in K^2
\]
と書け、$F$ は可逆である。$K^2$ はここでは二文字のテンソル積を表す。
$\ell(p)$ は $(p,p)$ が零点であることと $\nu p=p$ から両端で $p$ に消されるので、
$R$ を $\ell(p)$ に取れる。その二次行列は階数2である。実際 $w(p,-,-)$ の行列階数は、
末尾縮約 $r$ が同型であることから関係の $p$ での左評価階数である。
接平面での移動の微分は $D=F^{\mathsf t}$ である。零点近傍で混合二式から次点を解くと、
残る式の二次項は $u^{\mathsf t}R D u$ となる。三重点の重複度が3なのでこれは恒等零で、
$J=R D$ は可逆交代行列である。

$w$ の $z$ を一つ持つ部分を、左からは $zR+\sum A_{aj}x_a m_j$、
右からは $cRz+\sum B_{jb}m_j x_b$ の対応部分として比較する。
$z$ の先頭・中央・末尾の係数から $B=R$, $A=-FR$, $FRF^{\mathsf t}=cR$ を得る。
よって縮約を $R$ の係数で正規化して微分すると
$d\nu_p=cR^{-1}A^{\mathsf t}=-cR^{-1}R^{\mathsf t}D$ である。
$R=JD^{-1}$ と二次交代行列の恒等式から $c=\det D$、
$R^{-1}R^{\mathsf t}=-D^2/\det D$ なので $d\nu_p=D^3$ となる。
$d\sigma^4_p=D^4$ と比較して $D=1$ が従う。
特に三つの接方向は個別に保たれ、三成分の各々で移動は平行移動である。
三本の共点直線を個別に保ち、共点での微分が恒等の射影変換は、
共点 $[0:0:1]$ の座標で $[x:y:z]\mapsto[x:y:z+u x+v y]$ と書ける。
$\sigma^4$ のこの式の $u,v$ を4で割ると、各直線上の平行移動量の一意性から
$\sigma$ 自身を実現する。これも平面グラフになり矛盾する。
\end{proof}

\paragraph{横断二成分での特殊値と周期。}
上記 $h=1$ の九次元 $P_B$ のどの点でもAS正則である。辺の有限条件は、
二つの頂点で $d(p,p)\bmod kp$ がともに非零で、各成分の開トーラスで
$\det[p,\sigma p,d(p,\sigma p)]$ が恒等零でないことである。
実際、滑らかな部分では第四座標が二行から一意に解け、各根がその重複度のまま局所長になる。
頂点では節点基底と同じ検査になる。これらは一位置だけで検査できる。
一般の $p$ で $A_{\sigma^2p}(p)=B_p(\sigma^2p)\ne0$ であるから、
三次の混合条件は、$d_i(p,\sigma p)\in\langle p,\sigma p\rangle$ と
$d_{i+1}(\sigma p,\sigma^2p)\in\langle\sigma p,\sigma^2p\rangle$ を同値にする。
頂点でも剰余評価の非零性は移動で保たれる。
指定点の長さは20から両開トーラスの根の重複度を引いたものである。
有限部分は非空開である。両頂点で評価は接平面へ全射であり、
一般の補正では各枝の接線を避ける。各枝で移動の微分は1と異なるため、
その行列式には頂点の各逆像で単純な零があり、どの成分でも恒等零ではない。

同型は正規表示の射影補正を、二次曲線上の $s\mapsto l s$ および
$s\mapsto l/s$（移動の向きを対応させるもの）によって比較すれば必要十分である。
二次曲線と二つの交点を保つ射影変換はこれで尽きる。
元の座標で倍率の作用は対角 $(l,1,l^2)$ であり、補正への作用は
$d\mapsto GdG^{-1}$（値には二次の作用）で計算できる。
指定点以外の全局所環は曲線状なので、ここでも指定を忘れて同じ判定になる。
特定の係数での最小周期は $F_B^j$ に対してこの判定を行う。
倍率のトーラスは節点評価の二方向に異なる重みで作用するので、
射影空間の一般軌道は一次元、同型類の族は既約七次元である。

一般最小周期はやはり $m=2,4$ である。実際、上記の倍率の重み $-1$ の部分には、
$m=2$ で固有値 $1,-1$、$\lambda=\sqrt{-1}$ で $-1,-\sqrt{-1}$ の
$F_B$ 固有直線が各一本ある。重み3の部分には順に $1,-1$ および
$-1,\sqrt{-1}$ がある。これは前節の固有核計算を文字の重み $(1,0,2)$ で
分けて得られ、両重みに補生成元変更はない。
一つの重みに属する二係数の比は倍率で変わらず、移動の固有値比は位数 $m$ を持つ。
$m=2$ で自己同型に残る反転は重み $-1$ と3を交換し、
それらの係数比は一般には独立なので反転でも真の移動は戻せない。
以上の四種類は二つの楕円族とその境界の記述であり、非被約三次曲線をこの被覆には含めていない。
